\documentclass[11pt]{article}

\usepackage[margin=1in]{geometry}
\usepackage{listings}
\usepackage{xcolor}
\usepackage{hyperref}
\usepackage{booktabs}
\usepackage{amsmath}
\usepackage{enumitem}

\hypersetup{
    colorlinks=true,
    linkcolor=blue,
    urlcolor=blue,
    citecolor=blue
}

% ---- Listings style for C++ / Python ----
\definecolor{codegreen}{rgb}{0.0,0.5,0.0}
\definecolor{codegray}{rgb}{0.5,0.5,0.5}
\definecolor{codeblue}{rgb}{0.0,0.0,0.6}
\definecolor{backcolour}{rgb}{0.96,0.96,0.94}

\lstdefinestyle{ocv}{
    backgroundcolor=\color{backcolour},
    commentstyle=\color{codegreen},
    keywordstyle=\color{codeblue}\bfseries,
    numberstyle=\tiny\color{codegray},
    stringstyle=\color{codegreen},
    basicstyle=\ttfamily\footnotesize,
    breakatwhitespace=false,
    breaklines=true,
    captionpos=b,
    keepspaces=true,
    numbers=left,
    numbersep=6pt,
    showspaces=false,
    showstringspaces=false,
    showtabs=false,
    tabsize=4,
    frame=single,
    rulecolor=\color{codegray}
}
\lstset{style=ocv}

\title{\textbf{Holistically-Nested Edge Detection (HED) Module \\ for OpenCV \texttt{contrib}}\\[0.4em]
\large A Learning-Based Alternative to Canny Edge Detection}
\author{SASANK}
\date{June 2026}

\begin{document}
\maketitle

\begin{abstract}
This report documents a new OpenCV \texttt{contrib} module, \texttt{hed}, that adds
\emph{Holistically-Nested Edge Detection} (HED) as an alternative to the classical
Canny edge detector. Canny operates purely on local intensity gradients, so it
reacts to texture, noise, and lighting changes that do not correspond to real
object boundaries. HED instead uses a pretrained VGG-based convolutional network
that was trained on human-annotated boundaries, producing semantically meaningful
edge maps. The module wraps the network through OpenCV's \texttt{dnn} engine,
exposes a small \texttt{HEDDetector} API to both C++ and Python, and ships with a
demo, unit tests, and build integration.
\end{abstract}

\tableofcontents
\bigskip

%========================================================
\section{Motivation: Why Replace Canny?}
%========================================================

OpenCV's standard edge detector, \texttt{cv::Canny}, is a hand-crafted pipeline:
Gaussian smoothing, Sobel gradients, non-maximum suppression, and hysteresis
thresholding. It is fast and parameter-light, but it has well-known limitations:

\begin{itemize}[noitemsep]
    \item \textbf{Local-only decisions.} Each edge pixel is decided from a tiny
    neighbourhood. There is no notion of an \emph{object}; a strong texture edge
    and a true object boundary look identical.
    \item \textbf{Threshold sensitivity.} The two hysteresis thresholds must be
    tuned per image. Wrong values give either broken contours or a storm of
    spurious edges.
    \item \textbf{Noise and texture clutter.} Grass, fabric, foliage, and
    compression artifacts all produce gradients and therefore false edges.
    \item \textbf{No semantics.} Canny cannot tell that the outline of a person is
    more important than the wrinkles on their shirt.
\end{itemize}

HED addresses these by \emph{learning} what an edge is from data. It is a fully
convolutional network with deep supervision at multiple scales whose side-outputs
are fused into a single boundary probability map. Because it was trained on the
BSDS500 dataset of human-drawn boundaries, it produces clean, object-level edges
with no per-image threshold tuning.

%========================================================
\section{Summary of Improvements}
%========================================================

\begin{table}[h]
\centering
\begin{tabular}{@{}p{4.2cm}p{4.6cm}p{4.6cm}@{}}
\toprule
\textbf{Aspect} & \textbf{Canny (classic)} & \textbf{HED (this module)} \\
\midrule
Decision basis & Local intensity gradient & Learned multi-scale features \\
Semantics & None & Object-level boundaries \\
Parameters & 2 thresholds, tuned per image & None at inference \\
Noise / texture & Sensitive & Robust \\
Output & Binary edge mask & Continuous probability map $[0,1]$ \\
Scale awareness & Single scale & Multi-scale side outputs fused \\
Training data & Not applicable & Human-annotated boundaries (BSDS500) \\
\bottomrule
\end{tabular}
\caption{Canny vs.\ the HED module at a glance.}
\end{table}

\noindent\textbf{Key benefits in brief:}
\begin{enumerate}[noitemsep]
    \item \textbf{Semantic edges} — detects real object contours, not raw contrast.
    \item \textbf{No tuning} — no thresholds to set; the same model works across images.
    \item \textbf{Noise robustness} — texture and clutter are largely suppressed.
    \item \textbf{Soft confidence output} — a probability map ($CV\_32F$, $[0,1]$)
    that downstream code can threshold flexibly.
    \item \textbf{Seamless integration} — runs on OpenCV's existing \texttt{dnn}
    backend; no extra dependency, and it is exposed to Python automatically.
\end{enumerate}

%========================================================
\section{Module Layout}
%========================================================

The module follows the standard OpenCV \texttt{contrib} structure:

\begin{lstlisting}[language=bash, numbers=none]
modules/hed/
  CMakeLists.txt
  README.md
  include/opencv2/hed.hpp        # public API
  src/hed.cpp                    # implementation (dnn-backed)
  samples/hed_demo.py            # HED vs Canny demo
  test/test_hed.cpp              # unit tests
\end{lstlisting}

%========================================================
\section{Public API}
%========================================================

The public header declares a single abstract class, \texttt{HEDDetector}, with a
factory \texttt{create()} and one \texttt{detectEdges()} call. The
\texttt{CV\_EXPORTS\_W} / \texttt{CV\_WRAP} markers generate the Python bindings.

\begin{lstlisting}[language=C++, caption={include/opencv2/hed.hpp}]
#pragma once
#include "opencv2/core.hpp"
#include <string>

namespace cv {
namespace hed {

/** @defgroup hed Holistically-nested Edge Detection
This module implements HED edge detection as described in:

Xie, S., & Tu, Z. (2015). Holistically-nested edge detection.
Proceedings of the IEEE international conference on computer vision.
https://arxiv.org/abs/1504.06375
*/

//! @addtogroup hed
//! @{

/** @brief HED edge detector using a pretrained Caffe model.

Unlike classical detectors such as Canny, HED uses a VGG-based
convolutional network trained on human-annotated boundaries to
produce semantically meaningful edge maps.

@note Requires hed_pretrained_bsds.caffemodel and deploy.prototxt.
Download from: https://vcl.ucsd.edu/hed/
*/
class CV_EXPORTS_W HEDDetector {
public:
    /** @brief Creates an HEDDetector instance.
    @param modelPath path to hed_pretrained_bsds.caffemodel
    @param protoPath path to deploy.prototxt
    */
    CV_WRAP static cv::Ptr<HEDDetector> create(
        const std::string& modelPath,
        const std::string& protoPath
    );

    /** @brief Detects edges in an image.
    @param image input BGR image (CV_8UC3)
    @return edge map as CV_32F with values in [0, 1]
    */
    CV_WRAP virtual cv::Mat detectEdges(cv::InputArray image) = 0;

    /** @brief Destructor. */
    virtual ~HEDDetector() {}
};

//! @}

} // namespace hed
} // namespace cv
\end{lstlisting}

%========================================================
\section{Implementation}
%========================================================

The implementation loads the Caffe model through \texttt{cv::dnn::readNet} and runs
a forward pass. The input is converted to a blob with the network's mean-subtraction
values (the BGR means HED was trained with). The fused sigmoid output
(\texttt{"sigmoid-fuse"}) of shape $[1,1,H,W]$ is squeezed to an $H\times W$
single-channel float map and returned as an owned copy.

\begin{lstlisting}[language=C++, caption={src/hed.cpp}]
#include "opencv2/hed.hpp"
#include "opencv2/dnn.hpp"
#include "opencv2/imgproc.hpp"

namespace cv {
namespace hed {

class HEDDetectorImpl : public HEDDetector {
public:
    dnn::Net net;

    HEDDetectorImpl(const std::string& modelPath, const std::string& protoPath) {
        net = dnn::readNet(modelPath, protoPath);
    }

    cv::Mat detectEdges(cv::InputArray _image) override {
        cv::Mat image = _image.getMat();
        cv::Mat blob = dnn::blobFromImage(
            image, 1.0,
            cv::Size(image.cols, image.rows),
            cv::Scalar(104.00698793, 116.66876762, 122.67891434),
            false, false
        );
        net.setInput(blob);
        cv::Mat result = net.forward("sigmoid-fuse");

        // result shape is [1,1,H,W] -- squeeze to [H,W]
        cv::Mat edge(result.size[2], result.size[3], CV_32F, result.ptr<float>());
        cv::Mat output;
        edge.copyTo(output);
        return output;
    }
};

cv::Ptr<HEDDetector> HEDDetector::create(
    const std::string& modelPath,
    const std::string& protoPath)
{
    return cv::makePtr<HEDDetectorImpl>(modelPath, protoPath);
}

} // hed
} // cv
\end{lstlisting}

\paragraph{Implementation notes.}
\begin{itemize}[noitemsep]
    \item \textbf{Mean subtraction} uses the exact training means
    $(104.007,\,116.669,\,122.679)$, so the input statistics match what the
    network expects.
    \item \textbf{Full-resolution inference}: the blob is sized to the input image,
    so the edge map comes back at the original resolution.
    \item \textbf{Owned output}: \texttt{edge.copyTo(output)} copies out of the
    network's internal buffer so the returned \texttt{Mat} stays valid after the
    next forward pass.
\end{itemize}

%========================================================
\section{Build Integration}
%========================================================

The module declares its dependencies and enables Python wrapping in two lines.
The dependency on \texttt{opencv\_dnn} is what lets HED run without any external
framework.

\begin{lstlisting}[language=bash, caption={CMakeLists.txt}, numbers=none]
set(the_module opencv_hed)
ocv_define_module(hed opencv_core opencv_imgproc opencv_dnn WRAP python)
\end{lstlisting}

%========================================================
\section{Usage}
%========================================================

\subsection{C++}
\begin{lstlisting}[language=C++, numbers=none]
auto det = cv::hed::HEDDetector::create("model.caffemodel", "deploy.prototxt");
cv::Mat edges = det->detectEdges(image);   // CV_32F, values 0..1
\end{lstlisting}

\subsection{Python}
The bindings are generated automatically from the \texttt{CV\_WRAP} markers.

\begin{lstlisting}[language=Python, numbers=none]
detector = cv2.hed.HEDDetector.create("model.caffemodel", "deploy.prototxt")
edges = detector.detectEdges(image)        # float32, values 0..1
\end{lstlisting}

\subsection{Demo: HED vs.\ Canny}
The shipped sample runs both detectors on the same image for a side-by-side
comparison, making the qualitative improvement directly visible.

\begin{lstlisting}[language=Python, caption={samples/hed\_demo.py}]
import cv2
import sys
import argparse

parser = argparse.ArgumentParser()
parser.add_argument('--image', required=True)
parser.add_argument('--model', required=True)
parser.add_argument('--proto', required=True)
args = parser.parse_args()

image = cv2.imread(args.image)
if image is None:
    print(f"Cannot load image: {args.image}")
    sys.exit(1)

detector = cv2.hed.HEDDetector.create(args.model, args.proto)
edges = detector.detectEdges(image)
edges_u8 = (edges * 255).astype('uint8')

gray  = cv2.cvtColor(image, cv2.COLOR_BGR2GRAY)
canny = cv2.Canny(gray, 100, 200)

cv2.imshow("Original",        image)
cv2.imshow("Canny (classic)", canny)
cv2.imshow("HED (yours)",     edges_u8)
cv2.imwrite("hed_result.png", edges_u8)
print("Saved hed_result.png")
cv2.waitKey(0)
\end{lstlisting}

%========================================================
\section{Testing}
%========================================================

Two unit tests cover model loading and output contract. They skip gracefully when
the model files are not present in the test data path, so the suite stays green on
machines without the weights.

\begin{lstlisting}[language=C++, caption={test/test\_hed.cpp}]
#include "test_precomp.hpp"

namespace opencv_test { namespace {

TEST(HED, LoadModel)
{
    std::string model = cvtest::findDataFile("hed_pretrained_bsds.caffemodel", false);
    std::string proto = cvtest::findDataFile("deploy.prototxt", false);
    if (model.empty() || proto.empty())
        throw SkipTestException("HED model files not found");

    auto detector = cv::hed::HEDDetector::create(model, proto);
    ASSERT_FALSE(detector.empty());
}

TEST(HED, OutputShape)
{
    std::string model = cvtest::findDataFile("hed_pretrained_bsds.caffemodel", false);
    std::string proto = cvtest::findDataFile("deploy.prototxt", false);
    if (model.empty() || proto.empty())
        throw SkipTestException("HED model files not found");

    cv::Mat img = cv::Mat::zeros(100, 100, CV_8UC3);
    auto detector = cv::hed::HEDDetector::create(model, proto);
    cv::Mat edges = detector->detectEdges(img);

    EXPECT_EQ(edges.rows, 100);
    EXPECT_EQ(edges.cols, 100);
    EXPECT_EQ(edges.type(), CV_32F);
}

}} // namespace

CV_TEST_MAIN(".")
\end{lstlisting}

%========================================================
\section{Conclusion}
%========================================================

The \texttt{hed} module brings a modern, learning-based edge detector into the
OpenCV ecosystem as a drop-in alternative to Canny. By replacing hand-tuned
gradient thresholds with a network trained on human boundary annotations, it yields
edge maps that are cleaner, semantically meaningful, parameter-free at inference,
and robust to texture and noise. It integrates through the existing \texttt{dnn}
backend, ships C++ and auto-generated Python APIs, and includes a comparison demo
and unit tests --- making it ready for use in segmentation, contour extraction, and
other downstream pipelines where Canny's local heuristics fall short.

\paragraph{Reference.}
Xie, S., \& Tu, Z. (2015). \emph{Holistically-Nested Edge Detection}. ICCV.
\url{https://arxiv.org/abs/1504.06375}

\end{document}
